\section{Euler enclosures and the failure of an unrestricted constant one}
\label{rw:sec:euler}

Let $a,b>0$ with $w\ge1$. Define
\begin{equation}\label{rw:eq:eulerdef}
 f(n)=c_{2n+1}=\frac{H_{2n}^{(b)}}{(2n+1)^a},\qquad
 d_k=\sum_{n=0}^k(-1)^n\binom kn f(n),\qquad
 E_N=\sum_{k=0}^{N-1}\frac{d_k}{2^{k+1}}.
\end{equation}
The error is measured against $g_{a,b}=\Imag F_{a,b}(\iu)$.
On the critical line this is an Abel value; the ordinary alternating
series with terms $(-1)^nf(n)$ diverges by the term test.

\begin{theorem}[Extended one-sided Euler bounds]\label{rw:thm:euler}
For every $k\ge1$ one has $d_k<0$, while $d_0=0$. The $E_N$ decrease
strictly to $g_{a,b}$, and
\begin{equation}\label{rw:eq:exacteuler}
 E_N-g_{a,b}=-2^{-N}\int_{[0,1]}
       \frac{(1-u^2)^N}{1+u^2}\dd\nu_{a,b}(u)>0.
\end{equation}
Consequently
\begin{equation}\label{rw:eq:eulerenclosure}
 E_N-M(a,b)2^{-N}<g_{a,b}<E_N\qquad(N\ge1),
\end{equation}
with the following elementary choices:
\begin{equation}\label{rw:eq:Mchoices}
 M(a,b)=
 \begin{cases}
 1,&a\ge1,\\
 (1-b)^{-1},&0<a<1,\ 0<b<1,\ w\ge1,\\
 b/(b-1),&0<a<1,\ b>1,\\
 1+1/a,&0<a<1,\ b=1.
 \end{cases}
\end{equation}
On the critical line the sharper uniform choice is proved in
Corollary~\ref{subpick:cor:sharp-critical}; its matching subcritical
Euler theorem uses the positive difference measure.
The choice one in the first line is the prior optimal universal
constant on that restricted domain; it is not asserted to work in
the other lines.
\end{theorem}
\begin{proof}
The binomial theorem and \eqref{subunit:eq:moments} give
$d_k=\int(1-u^2)^k\dd\nu(u)$. Thus $d_0=0$, while the strictly
decreasing weight gives $d_k<0$ for $k\ge1$. With
$m=m(\nu)$, we have $|d_k|\le m$. Therefore the Euler series is
absolutely convergent and may be summed under the measure integral:
\[
 \sum_{k\ge0}\frac{d_k}{2^{k+1}}
 =\int\frac{\dd\nu(u)}{1+u^2}=\Imag F(\iu).
\]
Summing the geometric tail gives \eqref{rw:eq:exacteuler}. Its weight
is strictly decreasing, bounded by one, and strictly less than one
for every $u>0$. Since the negative Jordan part has no atom at zero,
$0<E_N-g<m2^{-N}$. It remains to bound $m$.

For $b<1,w>1$, rewrite \eqref{subunit:eq:generalK} as
\[
 K_{a,b}=\frac{k_{w-1}}{1-b}
          -\bigl[-\zeta(b)k_a+C_{a,b}\bigr].
\]
The two displayed densities are positive, and each has mass
$1/(1-b)$ in the $u$ measure, by zero total mass. Hence
$m<1/(1-b)$. On $w=1$ the exact Jordan mass is $1/a=1/(1-b)$.
For $b>1$, \eqref{subunit:eq:generalK} is a difference of positive densities
$\zeta(b)k_a$ and $k_{w-1}/(b-1)+C_{a,b}$, so
$m<\zeta(b)<1+1/(b-1)=b/(b-1)$.

For $b=1$ and $0<a<1$, use a translation estimate rather than the
singular limit of either of these bounds. Let $\sigma_a$ be the
probability law of $u=e^{-Y}$ with $Y\sim\mathrm{Gamma}(a,1)$.
Its logarithmic density $p_a(s)=e^{-s}s^{a-1}/\Gamma(a)$ decreases
strictly. If $D_{e^{-t}}$ denotes pushforward under $u\mapsto e^{-t}u$,
then direct comparison of $p_a(s)$ and $p_a(s-t)$ gives
\[
 \|\sigma_a-D_{e^{-t}}\sigma_a\|_{\mathrm{TV}}
 =2\int_0^tp_a(s)\dd s
 \le\frac{2t^a}{\Gamma(a+1)}.
\]
The signed-measure integral
\[
 \widetilde\nu=\int_0^\infty
       \frac{\sigma_a-D_{e^{-t}}\sigma_a}{e^t-1}\dd t
\]
converges in total variation. Its $(n-1)$st moment is
$n^{-a}\int_0^\infty(1-e^{-(n-1)t})(e^t-1)^{-1}\dd t=c_n$.
By uniqueness it is $\nu_{a,1}$, and
\[
 m\le\frac1{\Gamma(a+1)}\int_0^\infty\frac{t^a}{e^t-1}\dd t
 =\zeta(a+1)<1+1/a.
\]

Finally, for $a\ge1$ the earlier sharper bound can be recovered
without any fractional-order gap. At $a=1$ the equivalent base density is
\[
 K_{1,b}(T)=-\frac{T^b}{\Gamma(b+1)}
 +\frac1{\Gamma(b)}\int_T^\infty\frac{t^{b-1}}{e^t-1}\dd t.
\]
Its $n$th Laplace value is
$-n^{-b-1}+n^{-1}\sum_{j=1}^n j^{-b}=H_{n-1}^{(b)}/n$,
by Tonelli and a finite geometric sum. Moment uniqueness therefore
identifies this base density with \eqref{subunit:eq:generalK} or \eqref{subunit:eq:bOne}.
Both positive terms in this difference have $e^{-T}\dd T$ integral
one. For the second term, interchange the integrals and use
$(1-e^{-t})/(e^t-1)=e^{-t}$; the first is a gamma integral.
Their positive overlap makes the Jordan mass strictly less than one.
For $a>1$, convolution by $k_{a-1}$ in the $T$ coordinate preserves
moments after multiplication by $n^{-(a-1)}$ and contracts the
$L^1(e^{-T}\dd T)$ norm, because
$\int_0^\infty e^{-T}k_{a-1}(T)\dd T=1$.
Moment uniqueness identifies the result with the present kernel,
so $m<1$ for every real $a\ge1$. This proves all four choices.
\end{proof}

The same translation proof gives, when $0<a\le1$ and $a+b>1$,
the optional non-elementary bound
\[
 m\le\frac{\Gamma(a+b)\zeta(a+b)}{\Gamma(a+1)\Gamma(b)}.
\]
It may improve a particular estimate, but diverges as $a+b\downarrow1$.
The critical atomic construction is therefore not replaced by this
translation argument.

\begin{proposition}[Exact finite Euler weights]\label{rw:prop:weights}
For $N\ge1$,
\begin{equation}\label{rw:eq:weights}
 E_N=2^{-N}\sum_{n=0}^{N-1}(-1)^nf(n)
                  \sum_{j=n+1}^N\binom Nj.
\end{equation}
Thus the finite sum can be formed without constructing a triangular
difference table.
\end{proposition}
\begin{proof}
Substitute the definition of $d_k$ and interchange the finite sums.
The coefficient identity is
\[
 \sum_{k=n}^{N-1}2^{N-k-1}\binom kn=\sum_{j=n+1}^N\binom Nj,
\]
which follows by induction from Pascal's identity. Harmonic numbers
and binomial tails can then each be updated successively. This
rearrangement is already used in the integer-order manuscript; the
new feature in the implementation here is certified rational
interval arithmetic for fractional powers.
\end{proof}

\begin{theorem}[A certified scope obstruction]\label{rw:thm:counterexample}
For $a=1/10$ and $b=9/10$,
\begin{align}\label{rw:eq:counterinterval}
 -0.5297619322857712630946293655
 &<g_{1/10,9/10}\notag\\
 &<-0.5297619322857712630946293653.
\end{align}
In particular, $g_{1/10,9/10}<-13/25<-1/2$.
Since $E_1=0$, the inequality $E_N-2^{-N}<g_{a,b}$ fails at $N=1$
for this pair. A uniform constant one cannot be extended to the
entire region $a+b\ge1$. The failure also occurs strictly above
the threshold, with ordinary boundary convergence:
\begin{align}\label{rw:eq:counterordinary}
 -0.5308442529661593261720712602
 &<g_{1/10,1}\notag\\
 &<-0.5308442529661593261720712600<-0.53.
\end{align}
\end{theorem}
\begin{proof}
The accompanying standard-library program \texttt{code/certify.py}
evaluates \eqref{rw:eq:weights} with $N=96$ by outward rational interval
arithmetic. Every fractional power is bounded by integer tenth-root
inequalities, as detailed in Section~\ref{rw:sec:certificates}. The
critical mass bound is the exact rational $M=10$. Subtracting
$10\,2^{-96}$ from the lower endpoint of the finite Euler enclosure
and leaving its upper endpoint unchanged gives the interval recorded
in \texttt{data/exact\_certificates.json}. Outward decimal rounding
yields \eqref{rw:eq:counterinterval}; exact rational comparisons verify
its upper endpoint is less than $-13/25$. No numerical polylogarithm,
floating-point root, or unproved identity enters this certificate.
For $(a,b)=(1/10,1)$ the same computation uses $M=1+1/a=11$
and proves \eqref{rw:eq:counterordinary}. Its ordinary boundary
convergence follows from $a+b=11/10>1$ and
Corollary~\ref{rw:cor:boundary}.
\end{proof}

This is \emph{not} a counterexample to the baseline theorem, which
explicitly assumes $a\ge1$. It is a counterexample to dropping that
hypothesis while retaining the same universal error constant.

\begin{proposition}[The exact critical constant for fixed parameters]
\label{rw:prop:fixedconstant}
On $a+b=1$, put $d\sigma=-K_{1-b,b}(-\log u)\,du$. The scaled errors $2^N(E_N-g_{a,b})$ decrease strictly
to zero for $N\ge1$. The smallest constant in the non-strict estimate
$E_N-g_{a,b}\le C2^{-N}$, valid for all integers $N\ge1$ at that fixed
pair, is $C=-2g_{a,b}$.
\end{proposition}
\begin{proof}
The endpoint atom contributes zero to \eqref{rw:eq:exacteuler} when
$N\ge1$. Hence
\[
 2^N(E_N-g)=\int_0^1\frac{(1-u^2)^N}{1+u^2}\dd\sigma(u).
\]
Strict pointwise decrease and dominated convergence prove the first
claim. Its largest value occurs at $N=1$ and is $2(E_1-g)=-2g$.
For the corresponding strict inequality the infimum of admissible
constants is the same, but equality at $N=1$ prevents attaining it.
\end{proof}


\subsection{Exact arithmetic and independent native evidence}
\label{rw:sec:certificates}
The new standard-library verifier encloses every rational fractional power
by integer root inequalities on a $10^{100}$ grid. Products, harmonic sums
and finite Euler weights are then combined as exact rational intervals.
It certifies 616 strictly positive Hausdorff-difference signs, 112 negative
Euler increments and seven Gaussian analytic-value enclosures at 192
Euler terms. The six triangle cases use the rational budget $57/50$;
the remaining axis case uses $5/4$. The axis case $b=3/2$ also proves
$-2g_{0,3/2}>227/200$, supplying the finite arithmetic premise in the
axis-maximum theorem.

A separate native Wolfram calculation at 100 working digits evaluates
the outer Gamma integral of a depth-one polylogarithm, or the direct
axis value. All seven values lie inside their rational Euler enclosures.
The native quadrature is a numerical diagnostic, not interval quadrature.
The sign and remainder theorems depend on the written positive-measure
proofs. The original threshold report's tenth-root certificates are
preserved unchanged with its delivered code and data.
